%================================================================
\section{Interval geometry and the evaluation}\label{sec:evaluation}

\subsection{Two checks that answer different questions}\label{sec:geometry}

Each interval is checked in two ways, and the two checks answer different
questions. The first is \emph{containment}: does the interval contain the true
value of its target? The code records two flags. The flag \texttt{covers\_truth}
says whether the interval contains the truth at the point where the analyst works,
$\theta_2(c_0)$ or $\dspill(c_0)$. The flag \texttt{covers\_planted\_truth} says
whether it contains the truth at the real village, $\theta_2(\cstar)$ or
$\dspill(\cstar)$. At $\cstar$ the two coincide. Both flags describe one realized
panel. They are containment on one draw and are not estimates of
repeated-sampling coverage.

The second check is \emph{interval geometry}: how does the interval sit against
the image of the same functional? Intervals and images are subsets of the real
line, so a verdict is an inclusion. Let $J=[lo,hi]$ be an analyst's interval at
one point, let $\Theta=[L,U]$ be the image of its own target, and write
$\interior(\Theta)=(L,U)$. The function \texttt{classify\_interval} returns
\begin{align}
  \text{F-HONEST}
  &\iff lo\le L\ \text{and}\ hi\ge U
  \iff \Theta\subseteq J, \label{eq:honest}\\
  \text{F-OVERCONF}
  &\iff lo>L\ \text{and}\ hi<U
  \iff J\subset\interior(\Theta), \label{eq:overconf}
\end{align}
and no verdict in every other case, including an interval that overlaps the image
with neither inclusion. A0, A1 and A2 are compared with $\Theta_{\hm}$; Foil~2 is
compared with $\Theta_{\expset}$.

F-OVERCONF means that the interval is narrower than the set of values that the
targets leave open. It does not mean that the interval misses the truth. An
interval can contain the truth and be F-OVERCONF at the same time: it can be right
about its own world and still understate how much the targets leave open.
Conversely, an interval can miss the truth without being F-OVERCONF. We therefore
never use ``overconfident'' to mean ``misses the truth,'' and we report the two
checks side by side.

\subsection{Two researchers, two layers}\label{sec:twoanalysts}

The three diagnostic points serve two researchers who know different things.

\paragraph{The field experimenter} She has one panel from the real village, the
point $\cstar$, and she has calibrated nothing. Her question is whether her
interval contains the truth there. We call this layer~(a): is the estimator right
about its own world? The answer depends only on the identification conditions of
her estimator. The code also places her interval against the image, and we report
that comparison, but it is not a test of her estimator: she never used the model,
so a verdict against the image does not charge her with anything.

\paragraph{The calibrator} She fits the model to the targets and works at one
best-fitting point, the min-I point or the median-I point. Her interval is checked
three times. Layer~(a) asks whether it contains the truth at her own point. Layer~(b)
asks whether it contains the truth at the real village; this adds the maintained
assumption that her point is the world. The geometric verdict asks whether her
interval is narrower than what the targets leave open. She is the analyst who
replaces a set of not-ruled-out worlds with one of them, so the charge of
overconfidence applies to her.

\subsection{What the image maintains}

The image is honest only relative to what it maintains: the kernel, the box $C$,
the discrepancy $\delta$ and the targets with the threshold $\tau$. Change any of
them and the image changes. There is also a sampling caveat. The image is a range
over $\ThetaSample$ points drawn from a finite set of kept points, so it can
understate the range over the full level set of $I$. A verdict of F-OVERCONF
survives that error, because a larger image would still contain the interval
strictly inside. A verdict of F-HONEST would not survive it.

\subsection{Repeated panels for the three analysts of \texorpdfstring{$\theta_2$}{theta2}}\label{sec:repeated}

The flags above describe one panel. For A0, A1 and A2, amendment A2 adds a
repeated-sampling exercise at each diagnostic point: $R=\TradeoffR$ fresh panels
on a block of stream offsets that starts at $\TradeoffOffsetBase$ and is disjoint
from every other block in the project. Each panel redraws the simulation noise and,
for A1 and A2, the pilot's assignment. For each analyst the code reports the mean
estimate, the bias (mean minus $\theta_2(c_0)$), the standard deviation of the
estimates across panels (the sampling SD), the root mean squared error, the mean
reported standard error and its ratio to the sampling SD, the share of panels whose
interval contains $\theta_2(c_0)$ with an exact Clopper--Pearson $95\%$ interval
\citep{clopper1934}, the mean interval width, that width relative to the width of
$\Theta_{\hm}$, and, for A2, the number of undefined panels.

These quantities measure different things, and we keep them apart. The sampling
SD is the dispersion of the estimator. The mean reported SE is what the analyst
believes that dispersion to be, and the mean width is the interval she reports.
The repeated-sampling containment rate is the only coverage statement in this
part of the paper, and the truth it uses is $\theta_2(c_0)$, the truth at the
point that generated the panels. Addendum A2.1, written after these results were seen,
splits A1's and A2's reported variance into its parts; it is descriptive and is
reported with the results (Section~\ref{sec:a2results}).
